\chapter{The entropy connection and its boundary}\label{ch:entropy}

The canonical bridge can be extended to entropy in a precisely specified conservative thermal model. This chapter states that scope once, derives the signs and shows why Möbius contrasts inherit the same scaling. It also explains why a driven feedback service retains its finite EBU account without automatically acquiring this thermodynamic interpretation.


\section{A thermal statement needs a thermal model}
The word entropy is often used loosely for disorder, waste or difficulty. Here it has a precise role in a specified physical benchmark. A conservative overdamped system coupled to one reservoir at fixed temperature can exchange energy as heat while its stochastic system entropy changes. Under the accepted equilibrium conditions, both changes are related to the same endpoint potential difference, with opposite signs.

This is a narrow and useful connection. It does not say that every EBU improvement is universal entropy production or that an economy's total entropy vanishes. The earlier accounting identities remain valid wherever their own assumptions hold, even when this thermodynamic interpretation is unavailable.

\section{State the complete benchmark before the formula}
For each coalition entry, assume the same fixed conservative potential energy $U$, the same fixed temperature $T$, a conservative overdamped single-reservoir model, and the stationary canonical density at both endpoint times. Require compatible paths and boundary conventions, with no omitted external work, hidden reservoir, or additional entropy channel. Use $V=(U-U(x^*))/(k_BT)$ in one declared reference measure.

These assumptions must hold for every row of the coalition comparison. If applying a control changes the trap or injects work, the energy balance needs that work term. Merely knowing both endpoint potentials does not establish that all released energy went into the same medium in the way the benchmark assumes.

\begin{theorem}{Accepted P4 entropy relation}\label{thm:entropy}
Under the preceding conditions, for each compared process,
\begin{align}\label{eq:entropy-entrywise}
 \Delta s_{\mathrm{med}}(S)&=k_BE(S),\nonumber\\
 \Delta s_{\mathrm{sys}}(S)&=-k_BE(S),\\
 \Delta s_{\mathrm{tot}}(S)&=0.\nonumber
\end{align}
Consequently, applying the same Möbius transform to the entropy tables gives
\begin{align}\label{eq:entropy-mobius}
 m_{\Delta s_{\mathrm{med}}}(T)&=k_Bm_E(T),\nonumber\\
 m_{\Delta s_{\mathrm{sys}}}(T)&=-k_Bm_E(T),\\
 m_{\Delta s_{\mathrm{tot}}}(T)&=0.\nonumber
\end{align}
\end{theorem}
\begin{proof}
With no additional work channel, conservative heat delivered to the medium is $U(x_0)-U(x_S)$. Division by $T$ gives $k_BE(S)$. Stochastic system entropy is $s_{\mathrm{sys}}(x)=-k_B\ln p(x)$; its endpoint change is $-k_B\ln[p(x_S)/p(x_0)]=-k_BE(S)$. The two changes cancel. Möbius inversion is a linear signed sum, so applying it to each entrywise equality gives the coefficient relations.
\end{proof}

The second part is not a new thermodynamic law. It is linearity applied to an already conditional physical identity. Stochastic thermodynamics supplies the heat, work and entropy distinctions; the EBU notation organizes their relationship to the finite account \cite{seifert,maes}.

\section{Why the signs make sense}
Suppose the endpoint energy is lower. Under this benchmark, energy leaves the system as heat to the medium, so medium entropy increases and EBU is positive. The lower-energy endpoint has greater canonical density, so its stochastic surprisal and stochastic system entropy are lower. Their changes have opposite signs.

The cancellation does not mean that nothing happened, nor that every thermal trajectory in every model is reversible. It uses the stationary equilibrium density at both endpoint times and the specified conservative one-reservoir balance. A nonequilibrium initial distribution changes the stochastic system entropy calculation. Driving, multiple reservoirs or unrepresented internal variables can change the total balance.

A stochastic entropy change belongs to a specified realization; its ensemble average additionally uses the law of the process. The endpoint identity and its average therefore retain the physical model and probability convention that define them.

\section{An interaction changes units, not its logical meaning}
If an admissible pair in this benchmark has $m_E(AB)=-2$, its medium-entropy contrast is $-2k_B$. That is a contrast between the entropy changes of four declared processes. It is not necessarily the entropy change of one newly executed process called ``the interaction.'' The coefficient retains its counterfactual table meaning.

The same is true for a triple. Multiplying by $k_B$ gives the thermal units under the benchmark; it does not turn an eight-corner comparison into a causal allocation among three actors. The joint outcome must still be distinguished from an attribution or institutional payment.

\section{A current-carrying system exposes the limit}
A stationary density can agree with $e^{-V}$ while probability circulates. A driven ring can have uniform stationary probabilities and constant state potential while its clockwise and counterclockwise rate ratios carry nonzero cycle affinity. Its density-based EBU differences vanish, but its dynamics can produce entropy through the drive.

This is not a defect in the endpoint definition. It shows that a state-potential difference and a general nonequilibrium path functional are different objects. Appendix~\ref{app:probability} gives the ring explicitly. A loop returns the state potential to its start, while dissipative currents need not cancel their thermodynamic cost.

If a broader account includes nonconservative work, it must declare that work and its boundary. One cannot add an unspecified process charge to the fixed physical EBU definition and claim the original theorem proved the new expression. A separately defined net functional can be analyzed, but its additional terms require their own physical justification.

\section{Macrostate entropy is another declaration}
A constrained macrostate model may use a relation such as $S=S_{\mathrm{eq}}-\kappa V$. Its entropy can concern a coarse-grained ensemble with internal multiplicity. That object is not automatically the stochastic endpoint entropy $-k_B\ln p(x)$ used above. The ensemble, state resolution, reference measure and coefficient must be reconciled before equating them.

This distinction is important for EBU's broader ambition. A physically meaningful account of living systems may need several scales and internal states. The thermal benchmark offers a precise testbed; it does not license erasing those additional structures when moving to ecology or institutions.

\section{What the conditional connection is good for}
Within the stated P4 model, medium entropy changes and EBU values differ only by the factor $k_B$, while stochastic system entropy changes have the opposite sign. Linearity carries these relations through every coalition order. This is a complete theoretical consequence of the benchmark assumptions.

The distinction helps EBU keep physical claims precise. A thermal account can connect the finite unit to heat and temperature. A stock-and-pipeline controller can instead use its endpoint potential and stability function, each in its stated units. The tank's decaying oscillations supply no work--heat balance and no stationary canonical density, so they do not invoke P4.

The practical value is a modular account whose quantities retain their meanings. Heat, state improvement, accumulated exposure and controller storage can be displayed together where relevant without being added as though they were one substance. This supports decisions that recognise physical burdens at the appropriate level and makes the reason for each claimed connection clear.

We have completed the fixed-field chain. Two further theoretical questions remain before we assemble it: how to preserve the account when its yardstick changes, and how uncertainty in its inputs propagates through the exact formulas.
